\documentclass[11pt]{article}
\usepackage[a4paper,margin=0.65in]{geometry}
\usepackage[T1]{fontenc}
\usepackage{amsmath,amssymb,amsfonts}
\usepackage{graphicx}
\usepackage{pdfpages}
\usepackage[english,shorthands=off]{babel}
\usepackage{fancyhdr}
\usepackage[bookmarks=false,colorlinks=false]{hyperref}
\usepackage[protrusion=true,expansion=false]{microtype}
\emergencystretch=3em
\tolerance=600
\providecommand{\modd}{\mathrm{mod}}
\providecommand{\Disc}{\mathrm{Disc}}
\providecommand{\canont}{}
\providecommand{\asterism}{*}
\providecommand{\Hessian}{H}
\providecommand{\tome}{}
\providecommand{\page}{p.}
\pagestyle{fancy}\fancyhf{}\fancyhead[L]{Pages 526--545}\fancyhead[R]{Cayley --- Collected Papers, Vol.~I}\renewcommand{\headrulewidth}{0.4pt}
\setlength{\headheight}{15pt}

\begin{document}

\begin{center}
{\Large\bfseries The Collected Mathematical Papers of Arthur Cayley}\\[0.4em]
{\large Volume~I --- PDF pages 526--545 (book pp.~504--523)}\\[0.4em]
{\normalsize Papers [86] (cont.), [87], [88], [89], [90], [91]}
\end{center}

\vspace{1em}

% ===================================================================
% Paper [86] continuation -- ON THE DEVELOPABLE DERIVED FROM AN
% EQUATION OF THE FIFTH ORDER.  (concluding pages 504--506)
% ===================================================================

\section*{[86] On the Developable derived from an Equation of the Fifth Order \emph{(continued)}}

\noindent\textit{[Book p.~504]}\hfill

\smallskip

\noindent or the equations of the edge of regression are given by the system (equivalent of course to two equations),
\[
\left\| \begin{array}{ccccc} \alpha, & \beta, & \gamma, & A, & B \\ \beta, & \gamma, & \delta, & B, & C \end{array} \right\| = 0.
\]

The simplest mode of verifying \emph{\`a posteriori} that the edge of regression is only of the ninth order, appears to be to consider this curve as the common intersection of the three surfaces of the seventh order:
\[
\begin{aligned}
A^3 a - 3 A^2 B b + 3 A B^2 c - B^3 d &= 0,\\
A^3 b - 3 A^2 B c + 3 A B^2 d - B^3 e &= 0,\\
A^3 c - 3 A^2 B d + 3 A B^2 e - B^3 f &= 0,
\end{aligned}
\]
(which are at once obtained by combining the equation $B t + A = 0$ with the cubic equations in $t$). It is obvious from a preceding equation that if the equations first given are multiplied by $fA$, $-2eA + fB$, $2dA - eB$, and added, an identical result is obtained. This shows that the curve of the forty-ninth order, the intersection of the first two surfaces, is made up of the curve in question, the curve of the fourth order $A=0$, $B=0$ (which reckons for thirty-six, as being a triple line on each surface), and the curve which is common to the two surfaces of the seventh order and the surface $2 dA - eB = 0$. The equations of this last curve may be written,
\[
\begin{aligned}
e(af - 4 be + 3 cd) - 4 d (bf - 4 ce + 3 d^{2}) &= 0,\\
e^{2} a - 6 e^{2} d b + 12 e d^{2} c - 8 d^{4} &= 0,\\
e^{2} b - 6 e d^{2} c + 4 e d^{3} &= 0\,;
\end{aligned}
\]

or, observing that these equations are
\[
\begin{aligned}
f (a e - 4 b d) - 3 c (b e - 6 c e d + 4 d^{3}) &= 0,\\
e^{2}(a e - 4 b d) - 2 d (b e^{2} - 6 c e d + 4 d^{3}) &= 0,\\
e (b e^{2} - 6 c e d + 4 d^{3}) &= 0\,;
\end{aligned}
\]
the last-mentioned curve is the intersection of
\[
\begin{aligned}
a e - 4 b d &= 0,\\
b e^{2} - 6 c e d + 4 d^{3} &= 0,
\end{aligned}
\]
where the second surface contains the double line $e=0$, $d=0$, which is also a single line upon the first surface. Omitting this extraneous line, the intersection is of the fourth order; and we may remark that, in passing, it is determined (exclusively of the double line) as the intersection of the three surfaces
\[
\begin{aligned}
a e - 4 b d &= 0,\\
b e^{2} - 6 c e d + 4 d^{3} &= 0,\\
a^{2} d - 6 a b c + 4 b^{3} &= 0,
\end{aligned}
\]

\medskip
\noindent\textit{[Book p.~505]}\hfill

\smallskip

\noindent being in fact of the species iv.~4, of Mr Salmon's paper, ``On the Classification of Curves of Double Curvature'' [\emph{Journal}, vol.~v.\ (1850), pp.~23--46]. But returning to the question in hand; since $49 = 9 + 4 + 36$, we see that the curve common to the three surfaces of the seventh order, or the edge of regression, is, as it ought to be, of the ninth order. It only remains to express the equation of the developable surface in terms of the functions $A$, $B$, $C$, $\alpha$, $\beta$, $\gamma$, $\delta$, which determine the stationary points and edge of regression; I have satisfied myself that the required formula is
\[
\square = (AC - B^{2})^{2} - 1152\,\{A(\beta\delta - \gamma^{2}) - B(\alpha\delta - \beta\gamma) + C(\alpha\gamma - \beta^{2})\} = 0,
\]
where the quantities $\alpha$, $\beta$, $\gamma$, $\delta$ may be replaced by their values in $A$, $B$, $C$; and it will be noticed that when this is done, the terms of $\square$ are each of them at least of the third order in the last-mentioned functions.

I propose to term the family of developables treated of in this paper, `planar developables.' In general, the coefficients of the generating plane of a developable being algebraical functions of a variable parameter $t$, the equation rationalized with respect to the parameter belongs to a system of $n$ different planes; the developable which is the envelope of such a system may be termed a `multiplanar developable,' and in the particular case of $n$ being equal to unity, we have a planar developable. It would be very desirable to have some means of ascertaining from the equation of a developable what the degree of its `planarity' is.

\medskip

\noindent P.S.---At the time of writing the preceding paper I was under the impression that the only surface of the fourth order through the edge of regression was that given by the equation $AC - B^{2} = 0$; but Mr Salmon has since made known to me an entirely new form of the function $\square$, the component functions of which, equated to zero, give six different surfaces of the fourth order, each of them passing through the edge of regression. The form in question is
\[
3\,\square = L L' + 6 M M' - 6 N N',
\]
where
\[ \resizebox{\textwidth}{!}{$\displaystyle
\begin{aligned}
L  &= a^{2} f^{2} + 225 b^{2} e^{2} - 32 a c e^{2} - 32 b^{2} d f + 480 b c d e + 480 c^{3} e - 84 a b e f + 16 a c d f - 12 b d^{2} f - 12 a d^{2} e - 320 b c d e - 320 c^{2} d^{2},\\
L' &= 3 a^{2} f^{2} - 45 b^{2} e^{2} + 64 a c e^{2} + 64 b^{2} d f - 12 b c d e - 80 c^{3} e - 12 a b e f - 12 a c d f - 84 b c^{2} f - 8 a d^{2} e + 20 b c d e.
\end{aligned}
$} \]
\[ \resizebox{\textwidth}{!}{$\displaystyle
\begin{aligned}
M  &= 10 b c d^{2} - 18 a d^{3} - 15 b c^{2} e + 32 a c d e + 6 b^{2} d e - 9 a c^{2} f + 2 a b d f + a^{2} e f - 8 a b e^{2},\\
M' &= 10 c^{2} d e - 18 c f^{2} - 15 b d^{2} e + 22 b c d f + 6 a d^{3} - 9 b^{2} e f + 2 a c e f + a b f^{2} - 8 a d^{2} f,\\
N  &= 10 b^{3} d^{2} - 15 b^{2} c e - 12 a c d^{2} + 18 a c^{2} e + a b d e - 2 a^{2} b^{2} + 6 b^{2} f - 8 a b c f + 3 a^{2} d f,\\
N' &= 10 c^{2} e^{2} - 15 b d^{3} - 12 c^{2} d f + 18 b d^{2} f + b c e f - 2 b^{2} f^{2} + 6 a e^{2} - 8 a d e f + 3 a c f^{2},
\end{aligned}
$} \]
where it should be noticed that
\[
L + 3 L' = -10\,(A C - B^{2}).
\]

\medskip
\noindent\textit{[Book p.~506]}\hfill

\smallskip

\noindent The expressions of $L$, $L'$, $M$, $M'$, $N$, $N'$ as linear functions of $A$, $B$, $C$ (also due to Mr~Salmon) are
\[ \resizebox{\textwidth}{!}{$\displaystyle
\begin{aligned}
L  &= (11\,a e + 28\,b d - 39\,c^{2})\,A + (\,a f - 75\,b e + 74\,c d)\,B + (11\,b f + 28\,c e - 39\,d^{2})\,C,\\
L' &= (-7\,a e + 4\,b d + 3\,c^{2})\,A + (3\,a f + 10\,b e - 13\,c d)\,B + (-7\,b f + 4\,c e + 3\,d^{2})\,C,\\
M  &= 3\,(b c - a d)\,A + 3\,(a e - c^{2})\,B + (c d - a f)\,C,\\
M' &= (c d - a f)\,A + 3\,(b f - d^{2})\,B + 3\,(d e - c f)\,C,\\
N  &= 3\,(b^{2} - a c)\,A + 3\,(a d - b c)\,B + (b d - a e)\,C,\\
N' &= (c e - b f)\,A + 3\,(c f - d e)\,B + 3\,(e^{2} - d f)\,C.
\end{aligned}
$} \]

I propose resuming the subject of these forms, and the general theory, in a subsequent paper. [This was never written.]

\bigskip
\hrule
\bigskip

% ===================================================================
% Paper [87] -- NOTES ON ELLIPTIC FUNCTIONS (FROM JACOBI).
% Title page only (book p.~507).
% ===================================================================

\section*{[87] Notes on Elliptic Functions (from Jacobi)}

\noindent\textit{[Book p.~507]}\hfill

\smallskip

\noindent\emph{[From the Cambridge and Dublin Mathematical Journal, vol.~v. (1850), pp.~201--204.]}

\medskip

\noindent These Notes are mere translations from Jacobi's ``Note sur une nouvelle application de l'analyse des fonctions elliptiques \`a l'alg\`ebre,'' [\emph{Crelle}, t.~vii.\ (1831) pp.~41--43], and from the addition to the notice by him of the third supplement to Legendre's ``Th\'eorie des fonctions elliptiques'' [\emph{Crelle}, t.~viii.\ (1832) pp.~413--417].

\bigskip
\hrule
\bigskip

% ===================================================================
% Paper [88] -- ON THE TRANSFORMATION OF AN ELLIPTIC INTEGRAL.
% Book pp. 508--510.
% ===================================================================

\section*{[88] On the Transformation of an Elliptic Integral}

\noindent\textit{[Book p.~508]}\hfill

\smallskip

\noindent\emph{[From the Cambridge and Dublin Mathematical Journal, vol.~v. (1850), pp.~204--206.]}

\medskip

\noindent The following is a demonstration of a formula proved incidentally by Mr~Boole (\emph{Journal}, vol.~ii.\ [1847] p.~7), in a paper ``On the Attraction of a Solid of Revolution on an External Point.''

\medskip

Let
\[
U = \int_{-1}^{+1} \frac{dx}{\sqrt{(1-x^{2})\{1-(m x + n)^{2}\}}}\,;
\]
then, assuming
\[
x = \frac{\alpha + i y}{1 - i \alpha y}\,,
\]
(so that $x = \pm 1$ gives $y = \pm 1$), we obtain
\[
1 - x^{2} = \frac{(1 + \alpha^{2})(1 - y^{2})}{(1 - i \alpha y)^{2}}\,,
\]
\[
m x + n = \frac{(n - i m \alpha) + (m - i n \alpha)\,y}{1 - i \alpha y}\,.
\]
Assume therefore
\[
i \alpha + (n - i m \alpha)(m - i n \alpha) = 0,
\]
whence
\[
-i \alpha = \frac{(1 - m^{2} - n^{2}) + \Delta}{2 m n}\qquad (\Delta^{2} = 1 + m^{4} + n^{4} - 2 m^{2} - 2 n^{2} - 2 m^{2} n^{2}).
\]

\medskip
\noindent\textit{[Book p.~509]}\hfill

\smallskip

\noindent we find
\[
1 - (m x + n)^{2} = \frac{1 - (n - i m \alpha)^{2}}{(1 - i \alpha y)^{2}}\,\bigl\{1 - (m - i n \alpha)^{2} y^{2}\bigr\},
\]
and also
\[
dx = \frac{(1 + \alpha^{2})\,dy}{(1 - i \alpha y)^{2}}\,,
\]
whence
\[
U = \sqrt{\!\left[\frac{1 + \alpha^{2}}{1 - (n - i m \alpha)^{2}}\right]}\;\int_{-1}^{+1} \frac{dy}{\sqrt{(1 - y^{2})\{1 - (m - i n \alpha)^{2} y^{2}\}}}\,,
\]
that is
\[
U = 2 \sqrt{\!\left[\frac{1 + \alpha^{2}}{1 - (n - i m \alpha)^{2}}\right]}\;\int_{0}^{1} \frac{dy}{\sqrt{(1 - y^{2})\{1 - (m - i n \alpha)^{2} y^{2}\}}}\,.
\]
But since
\[
n - i m \alpha = \frac{1 - m^{2} + n^{2} + \Delta}{2 n}\,,
\]
\[
m - i n \alpha = \frac{1 + m^{2} - n^{2} + \Delta}{2 m}\,,
\]
we have
\[
1 - (n - i m \alpha)^{2} = -\,\frac{\Delta}{2 n^{2}}\,(\Delta + 1 - m^{2} + n^{2}),
\]
\[
1 + \alpha^{2} = -\,\frac{\Delta}{2 m^{2} n^{2}}\,(\Delta + 1 - m^{2} - n^{2})\,;
\]
and therefore
\[ \resizebox{\textwidth}{!}{$\displaystyle
\frac{1 + \alpha^{2}}{1 - (n - i m \alpha)^{2}} = \frac{1}{m^{2}}\cdot\frac{\Delta + 1 - m^{2} - n^{2}}{\Delta + 1 - m^{2} + n^{2}}
= \frac{1}{m^{2}}\cdot\frac{(1 - m^{2} - n^{2} + \Delta)(1 - m^{2} + n^{2} - \Delta)}{(1 - m^{2} + n^{2} + \Delta)(1 - m^{2} + n^{2} - \Delta)} = \frac{2(1 + m^{2} - n^{2} + \Delta)}{4 m^{2}}\,;
$} \]
consequently
\[
U = \frac{1}{m}\sqrt{\{2(1 + m^{2} - n^{2} + \Delta)\}}\;\int_{0}^{1} \frac{dy}{\sqrt{\Bigl[(1 - y^{2})\bigl\{1 - \dfrac{1 + m^{2} - n^{2} + \Delta}{2 m}\, y^{2}\bigr\}\Bigr]}}\,.
\]

Write
\[
\frac{1}{\lambda} = \frac{1 + m^{2} - n^{2} + \Delta}{2 m}\,,\qquad \lambda^{2} k^{2} = \frac{4 m n}{(1 + m)^{2} - n^{2}}\,;
\]
then
\[
U = \frac{\sqrt{4\lambda k}}{\lambda}\cdot\frac{1}{\sqrt{\{(1 + m)^{2} - n^{2}\}}}\;\int_{0}^{1} \frac{dy}{\sqrt{(1 - y^{2})(1 - \lambda^{2} y^{2})}}\,,
\]
where $\lambda$ and $k$ are connected by the relation that exists for the transformation of the second order, viz.\
\[
\lambda = \frac{2\sqrt{k}}{1 + k}\,.
\]

\medskip
\noindent\textit{[Book p.~510]}\hfill

\smallskip

\noindent as may be immediately verified: hence, assuming
\[
y = \frac{\sqrt{\lambda}\, z}{\sqrt{k}\,\sqrt{(1 - \lambda^{2} z^{2})}}\,,
\]
which gives
\[
1 - y^{2} = \frac{1 - z^{2}}{\sqrt{(1 - \lambda^{2} z^{2})}}\,,
\]
we find
\[
U = \frac{4}{\sqrt{\{(1 + m)^{2} - n^{2}\}}}\;\int_{0}^{1} \frac{dz}{\sqrt{\Bigl[(1 - z^{2})\Bigl\{1 - \dfrac{4 m n}{(1 + m)^{2} - n^{2}}\, z^{2}\Bigr\}\Bigr]}}\,,
\]
that is
\[
\int_{-1}^{+1} \frac{dx}{\sqrt{(1 - x^{2})\{1 - (m x + n)^{2}\}}} = 4 \int_{0}^{1} \frac{dz}{\sqrt{(1 - z^{2})\{(1 + m)^{2} - n^{2} - 4 m n\, z^{2}\}}}\,.
\]
Writing here
\[
x = \cos\theta,\qquad z = \cos\tfrac{1}{2}\phi,
\]
then
\[
\int_{0}^{\pi} \frac{d\theta}{\sqrt{1 - (m\cos\theta + n)^{2}}} = \int_{0}^{\pi} \frac{d\phi}{\sqrt{1 + m^{2} - n^{2} - 2 m n \cos\phi}}\,,
\]
or if
\[
m = \frac{r}{a}\,,\qquad n = \frac{z}{a}\,,
\]
then finally
\[
\int_{0}^{\pi} \frac{d\theta}{\sqrt{\{a^{2} + (z + r\cos\theta)^{2}\}}} = \int_{0}^{\pi} \frac{d\phi}{\sqrt{(a^{2} + r^{2} + z^{2} - 2 a r \cos\phi)}}\,,
\]
the formula in question.

\bigskip
\hrule
\bigskip

% ===================================================================
% Paper [89] -- ON THE ATTRACTION OF ELLIPSOIDS (JACOBI'S METHOD).
% Book pp. 511--518.
% ===================================================================

\section*{[89] On the Attraction of Ellipsoids (Jacobi's Method)}

\noindent\textit{[Book p.~511]}\hfill

\smallskip

\noindent\emph{[From the Cambridge and Dublin Mathematical Journal, vol.~v. (1850), pp.~217--226.]}

\medskip

\noindent In a letter published in 1846 in Liouville's \emph{Journal} (t.~xi.\ p.~341) Jacobi says, ``Il y a quatorze ans, je me suis pos\'e le probl\`eme de chercher l'attraction d'un ellipso\"ide homog\`ene exerc\'ee sur un point ext\'erieur quelconque par une m\'ethode analogue \`a celle employ\'ee par Maclaurin par rapport aux points situ\'es dans les axes principaux. J'y suis parvenu par trois substitutions cons\'ecutives. La premi\`ere est une transformation de coordonn\'ees; par la seconde le radical $\sqrt{(1 - m^{2} \sin^{2}\beta \cos^{2}\psi - n^{2} \sin^{2}\beta \sin^{2}\psi)}$ qui entre dans la double int\'egrale transform\'ee est rendu rationnel au moyen de la double substitution
\[
m \sin\beta \cos\psi = \sin v \cos\theta,\qquad m \sin\beta \sin\psi = \sin v \sin\theta\,;
\]
la troisi\`eme est encore une transformation de coordonn\'ees. La recherche du sens g\'eom\'etrique de ces trois substitutions m'a conduit \`a approfondir la th\'eorie des surfaces confocales par rapport auxquelles je d\'ecouvris quantit\'e de beaux th\'eor\`emes dont je communiquai quelques-uns des principaux \`a M.~Steiner. Consid\'erons l'ellipso\"ide confocal men\'e par le point attir\'e $P$ et le point $p$ de l'ellipso\"ide propos\'e, conjugu\'e \`a $P$. Soient $Q$ et $q$ deux autres points conjugu\'es quelconques situ\'es respectivement sur l'ellipso\"ide ext\'erieur et int\'erieur. Menons de $P$ un premier c\^one tangent \`a l'ellipso\"ide int\'erieur, de $p$ un second c\^one tangent \`a l'ellipso\"ide ext\'erieur. Ce dernier, tout imaginaire qu'il est, a ses trois axes r\'eels (ainsi que ses deux droites focales). La premi\`ere substitution ram\`ene les axes de l'ellipso\"ide \`a ceux du premier c\^one (c'est la substitution employ\'ee par Poisson, mais que j'avais ant\'erieurement trait\'ee et m\^eme \'etendue \`a un nombre quelconque de variables dans le m\'emoire \emph{De binis Functionibus homogeneis \&c.} [\emph{Crelle}, t.~xii.\ (1834) pp.~1--69]). Par la seconde substitution les angles que la droite $Pq$ forme avec les axes du premier c\^one sont ramen\'es aux angles que la droite $pQ$ forme avec les axes du second. Par la derni\`ere substitution, on retourne de ces axes aux axes de l'ellipso\"ide. La seconde substitution r\'epond \`a un th\'eor\`eme de g\'eom\'etrie remarquable, savoir que `Les cosinus des angles que la droite $Pq$ forme avec deux des axes du premier c\^one sont en raison constante avec les cosinus des angles que la droite

\medskip
\noindent\textit{[Book p.~512]}\hfill

\smallskip

\noindent $pQ$ forme avec deux des axes du second c\^one; ces deux axes sont les tangentes situ\'es respectivement dans les sections de plus grande et de moindre courbure de chaque ellipso\"ide, le troisi\`eme axe \'etant la normale \`a l'ellipso\"ide.' Tout cela semble difficile \`a \'etablir par la synth\`ese.''

\medskip

The object of this paper is to develope the above method of finding the attraction of an ellipsoid.

\medskip

Consider an exterior ellipsoid, the squared semiaxes of which are $f + u$, $g + u$, $h + u$; and an interior ellipsoid, the squared semiaxes of which are $f + \bar{u}$, $g + \bar{u}$, $h + \bar{u}$. Let $u$, $p$, $q$ be the elliptic coordinates of a point $P$ on the exterior ellipsoid, the elliptic coordinates of the corresponding point $\bar{P}$ on the interior ellipsoid will be $\bar{u}$, $p$, $q$, and if $a$, $b$, $c$ and $\bar{a}$, $\bar{b}$, $\bar{c}$ represent the ordinary coordinates of these points (the principal axes being the axes of coordinates), we have
\[
\begin{aligned}
a^{2} &= \frac{(f + u)(f + q)(f + r)}{(f - g)(f - h)}, & \bar{a}^{2} &= \frac{(f + \bar{u})(f + q)(f + r)}{(f - g)(f - h)}\,,\\[2pt]
b^{2} &= \frac{(g + u)(g + q)(g + r)}{(g - h)(g - f)}, & \bar{b}^{2} &= \frac{(g + \bar{u})(g + q)(g + r)}{(g - h)(g - f)}\,,\\[2pt]
c^{2} &= \frac{(h + u)(h + q)(h + r)}{(h - f)(h - g)}, & \bar{c}^{2} &= \frac{(h + \bar{u})(h + q)(h + r)}{(h - f)(h - g)}\,.
\end{aligned}
\]

I form the systems of equations
\[
\begin{aligned}
a_{1}^{2} &= \frac{(u + f)(u + g)(u + h)}{(u - q)(u - r)}, & \bar{a}_{1}^{2} &= \frac{(u + f)(u + g)(u + h)}{(u - q)(u - r)}\,,\\[2pt]
b_{1}^{2} &= \frac{(q + f)(q + g)(q + h)}{(q - r)(q - u)}, & \bar{b}_{1}^{2} &= \frac{(q + f)(q + g)(q + h)}{(q - r)(q - \bar{u})}\,,\\[2pt]
c_{1}^{2} &= \frac{(r + f)(r + g)(r + h)}{(r - u)(r - q)}, & \bar{c}_{1}^{2} &= \frac{(r + f)(r + g)(r + h)}{(r - \bar{u})(r - q)}\,.
\end{aligned}
\]

\[
\alpha = \frac{a_{1} a}{f + u},\qquad \beta = \frac{a_{1} b}{g + u},\qquad \gamma = \frac{a_{1} c}{h + u},
\]
\[
\alpha' = \frac{b_{1} a}{f + q},\qquad \beta' = \frac{b_{1} b}{g + q},\qquad \gamma' = \frac{b_{1} c}{h + q},
\]
\[
\alpha'' = \frac{c_{1} a}{f + r},\qquad \beta'' = \frac{c_{1} b}{g + r},\qquad \gamma'' = \frac{c_{1} c}{h + r},
\]
\[
\bar{\alpha} = \frac{\bar{a}_{1} \bar{a}}{f + \bar{u}},\qquad \bar{\beta} = \frac{\bar{a}_{1} \bar{b}}{g + \bar{u}},\qquad \bar{\gamma} = \frac{\bar{a}_{1} \bar{c}}{h + \bar{u}},
\]
\[
\bar{\alpha}' = \frac{\bar{b}_{1} \bar{a}}{f + q},\qquad \bar{\beta}' = \frac{\bar{b}_{1} \bar{b}}{g + q},\qquad \bar{\gamma}' = \frac{\bar{b}_{1} \bar{c}}{h + q},
\]
\[
\bar{\alpha}'' = \frac{\bar{c}_{1} \bar{a}}{f + r},\qquad \bar{\beta}'' = \frac{\bar{c}_{1} \bar{b}}{g + r},\qquad \bar{\gamma}'' = \frac{\bar{c}_{1} \bar{c}}{h + r}.
\]

\medskip
\noindent\textit{[Book p.~513]}\hfill

\smallskip

\noindent And then writing
\[
\begin{aligned}
X &= \alpha X_{1} + \alpha' Y_{1} + \alpha'' Z_{1}, & \bar{X} &= \bar{\alpha}\,\bar{X}_{1} + \bar{\alpha}'\,\bar{Y}_{1} + \bar{\alpha}''\,\bar{Z}_{1},\\
Y &= \beta X_{1} + \beta' Y_{1} + \beta'' Z_{1}, & \bar{Y} &= \bar{\beta}\,\bar{X}_{1} + \bar{\beta}'\,\bar{Y}_{1} + \bar{\beta}''\,\bar{Z}_{1},\\
Z &= \gamma X_{1} + \gamma' Y_{1} + \gamma'' Z_{1}, & \bar{Z} &= \bar{\gamma}\,\bar{X}_{1} + \bar{\gamma}'\,\bar{Y}_{1} + \bar{\gamma}''\,\bar{Z}_{1},
\end{aligned}
\]
if $X$, $Y$, $Z$ are the cosines of the inclinations of a line $PQ$ to the principal axes of the ellipsoids, $X_{1}$, $Y_{1}$, $Z_{1}$ will be the cosines of the inclinations of this line to the principal axes of the cone having $P$ for its vertex, and circumscribed about the interior ellipsoid. In like manner, $\bar{X}$, $\bar{Y}$, $\bar{Z}$ being the cosines of the inclinations of a line $\bar{P}\bar{Q}$ to the principal axes of the ellipsoids, $\bar{X}_{1}$, $\bar{Y}_{1}$, $\bar{Z}_{1}$ will be the cosines of the inclinations of this line to the principal axes of the cone having $\bar{P}$ for its vertex and circumscribed about the exterior ellipsoid. Assuming that the points $Q$, $\bar{Q}$ are situated upon the exterior and interior ellipsoids respectively, suppose that $X_{1}$, $Y_{1}$, $Z_{1}$ and $\bar{X}_{1}$, $\bar{Y}_{1}$, $\bar{Z}_{1}$ are connected by the equivalent systems of equations,
\[
X_{1} = \sqrt{(u - \bar{u})}\,\sqrt{\!\left(\frac{\bar{X}_{1}^{\,2}}{u - \bar{u}} + \frac{\bar{Y}_{1}^{\,2}}{\bar{u} - q} + \frac{\bar{Z}_{1}^{\,2}}{\bar{u} - r}\right)},
\]
\[
Y_{1} = \sqrt{\!\left(\frac{\bar{u} - q}{u - q}\right)}\,\bar{Y}_{1},\qquad Z_{1} = \sqrt{\!\left(\frac{\bar{u} - r}{u - r}\right)}\,\bar{Z}_{1},
\]
\[
\bar{X}_{1} = \sqrt{(u - \bar{u})}\,\sqrt{\!\left(\frac{X_{1}^{\,2}}{u - \bar{u}} + \frac{Y_{1}^{\,2}}{u - q} + \frac{Z_{1}^{\,2}}{u - r}\right)},
\]
\[
\bar{Y}_{1} = \sqrt{\!\left(\frac{u - q}{\bar{u} - q}\right)}\,Y_{1},\qquad \bar{Z}_{1} = \sqrt{\!\left(\frac{u - r}{\bar{u} - r}\right)}\,Z_{1},
\]
then it will presently be shown that the points $Q$, $\bar{Q}$ are corresponding points, which will prove the geometrical theorem of Jacobi. Before proceeding further it will be convenient to notice the formulae
\[
1 - \frac{a^{2}}{f + u} - \frac{b^{2}}{g + u} - \frac{c^{2}}{h + u} = \frac{\bar{u} - u}{a_{1}^{\,2}}\,,
\]
\[
\frac{Xa}{f + u} + \frac{Yb}{g + u} + \frac{Zc}{h + u} = \frac{\bar{u} - u}{a_{1}^{\,2}}\!\left(\frac{X_{1} a_{1}}{u - \bar{u}} + \frac{Y_{1} b_{1}}{u - q} + \frac{Z_{1} c_{1}}{u - r}\right),
\]
\[
\left(\frac{Xa}{f + u} + \frac{Yb}{g + u} + \frac{Zc}{h + u}\right)^{\!2} + \frac{\bar{u} - u}{a_{1}^{\,2}}\!\left(\frac{X^{2}}{f + u} + \frac{Y^{2}}{g + u} + \frac{Z^{2}}{h + u}\right) = \frac{\bar{u} - u}{a_{1}^{\,2}}\!\left(\frac{X_{1}^{2}}{u - \bar{u}} + \frac{Y_{1}^{2}}{u - q} + \frac{Z_{1}^{2}}{u - r}\right).
\]

\medskip
\noindent\textit{[Book p.~514]}\hfill

\smallskip

\noindent and the corresponding ones
\[
1 - \frac{\bar{a}^{2}}{f + \bar{u}} - \frac{\bar{b}^{2}}{g + \bar{u}} - \frac{\bar{c}^{2}}{h + \bar{u}} = \frac{u - \bar{u}}{\bar{a}_{1}^{\,2}}\,,
\]
\[
\frac{\bar{X}\bar{a}}{f + \bar{u}} + \frac{\bar{Y}\bar{b}}{g + \bar{u}} + \frac{\bar{Z}\bar{c}}{h + \bar{u}} = \frac{u - \bar{u}}{\bar{a}_{1}^{\,2}}\!\left(\frac{\bar{X}_{1} \bar{a}_{1}}{\bar{u} - u} + \frac{\bar{Y}_{1} \bar{b}_{1}}{\bar{u} - q} + \frac{\bar{Z}_{1} \bar{c}_{1}}{\bar{u} - r}\right),
\]
\[
\left(\frac{\bar{X}\bar{a}}{f + \bar{u}} + \frac{\bar{Y}\bar{b}}{g + \bar{u}} + \frac{\bar{Z}\bar{c}}{h + \bar{u}}\right)^{\!2} + \frac{u - \bar{u}}{\bar{a}_{1}^{\,2}}\!\left(\frac{\bar{X}^{2}}{f + \bar{u}} + \frac{\bar{Y}^{2}}{g + \bar{u}} + \frac{\bar{Z}^{2}}{h + \bar{u}}\right) = \frac{u - \bar{u}}{\bar{a}_{1}^{\,2}}\!\left(\frac{\bar{X}_{1}^{2}}{\bar{u} - u} + \frac{\bar{Y}_{1}^{2}}{\bar{u} - q} + \frac{\bar{Z}_{1}^{2}}{\bar{u} - r}\right).
\]

The coordinates of the point $Q$ are obviously $a + \rho X$, $b + \rho Y$, $c + \rho Z$ (where $\rho = PQ$); substituting these values in the equation of the interior ellipsoid, we obtain
\[
\rho\,\!\left\{\frac{(u - \bar{u})}{a_{1}^{\,2}} + 2\rho\!\left(\frac{Xa}{f + \bar{u}} + \frac{Yb}{g + \bar{u}} + \frac{Zc}{h + \bar{u}}\right) + \rho^{2}\!\left(\frac{X^{2}}{f + \bar{u}} + \frac{Y^{2}}{g + \bar{u}} + \frac{Z^{2}}{h + \bar{u}} - 1\right) = 0\right\};
\]
reducing the coefficients of this equation by the formulae first given, and omitting a factor $\dfrac{\bar{u} - u}{\bar{a}_{1}^{\,2}}$, we obtain
\[
\frac{a_{1}^{2} \bar{X}_{1}^{2}}{(\bar{u} - u)^{2}} - \!\left(\frac{X_{1} a_{1}}{\bar{u} - u} - \frac{Y_{1} b_{1}}{\bar{u} - q} + \frac{Z_{1} c_{1}}{\bar{u} - r}\right)^{\!2} \rho^{2} + 2 \rho\!\left(\frac{X_{1} a_{1}}{\bar{u} - u} + \frac{Y_{1} b_{1}}{\bar{u} - q} + \frac{Z_{1} c_{1}}{\bar{u} - r}\right) - 1 = 0,
\]
that is,
\[
\frac{a_{1}^{2} \bar{X}_{1}^{2}}{(\bar{u} - u)^{2}}\,\rho^{2} = \!\left\{\rho\!\left(\frac{X_{1} a_{1}}{\bar{u} - u} - \frac{Y_{1} b_{1}}{\bar{u} - q} + \frac{Z_{1} c_{1}}{\bar{u} - r}\right) - 1\right\}^{\!2},
\]
or
\[
\rho = \frac{1}{\dfrac{X_{1} a_{1} - X_{1}\bar{a}_{1}}{\bar{u} - u} + \dfrac{Y_{1} b_{1}}{\bar{u} - q} + \dfrac{Z_{1} c_{1}}{\bar{u} - r}}\,,
\]
which is easily transformed into
\[
\rho = \frac{\bar{u} - u}{a_{1} X_{1} - \bar{a}_{1} \bar{X}_{1} + \dfrac{(f + u) b_{1} Y_{1} - (f + \bar{u}) \bar{b}_{1} \bar{Y}_{1}}{f + q} + \dfrac{(f + u) c_{1} Z_{1} - (f + \bar{u}) \bar{c}_{1} \bar{Z}_{1}}{f + r}}\,;
\]
and this form remaining unaltered when $u$ and $\bar{u}$ are interchanged, it follows that if $PQ = \rho$, then $\rho = \bar{\rho}$, which is a known theorem. The value of $\rho$ or $\bar{\rho}$ may however be expressed in a yet simpler form; for, considering the expression
\[ \resizebox{\textwidth}{!}{$\displaystyle
\frac{X}{\sqrt{(f + u)}} - \frac{\bar{X}}{\sqrt{(f + \bar{u})}}
= \frac{a}{\sqrt{(f + u)}}\!\left\{\frac{a_{1} X_{1}}{f + u} + \frac{b_{1} Y_{1}}{f + q} + \frac{c_{1} Z_{1}}{f + r}\right\}
- \frac{\bar{a}}{\sqrt{(f + \bar{u})}}\!\left\{\frac{\bar{a}_{1} \bar{X}_{1}}{f + \bar{u}} + \frac{\bar{b}_{1} \bar{Y}_{1}}{f + q} + \frac{\bar{c}_{1} \bar{Z}_{1}}{f + r}\right\}
$} \]
\[
= \frac{-1}{\bar{u} - u}\!\left\{\frac{a}{\sqrt{(f + u)}} - \frac{\bar{a}}{\sqrt{(f + \bar{u})}}\right\}
\]
\[ \resizebox{\textwidth}{!}{$\displaystyle
\times \!\left\{a_{1} X_{1} - \bar{a}_{1} \bar{X}_{1} + \frac{(f + u) b_{1} Y_{1} - (f + \bar{u}) \bar{b}_{1} \bar{Y}_{1}}{f + q} + \frac{(f + u) c_{1} Z_{1} - (f + \bar{u}) \bar{c}_{1} \bar{Z}_{1}}{f + r}\right\}\,;
$} \]

\medskip
\noindent\textit{[Book p.~515]}\hfill

\smallskip

\noindent we see that
\[
\frac{X}{\sqrt{(f + u)}} - \frac{\bar{X}}{\sqrt{(f + \bar{u})}} = -\,\frac{1}{\rho}\!\left\{\frac{a}{\sqrt{(f + u)}} - \frac{\bar{a}}{\sqrt{(f + \bar{u})}}\right\},
\]
and similarly
\[
\frac{Y}{\sqrt{(g + u)}} - \frac{\bar{Y}}{\sqrt{(g + \bar{u})}} = -\,\frac{1}{\rho}\!\left\{\frac{b}{\sqrt{(g + u)}} - \frac{\bar{b}}{\sqrt{(g + \bar{u})}}\right\},
\]
\[
\frac{Z}{\sqrt{(h + u)}} - \frac{\bar{Z}}{\sqrt{(h + \bar{u})}} = -\,\frac{1}{\rho}\!\left\{\frac{c}{\sqrt{(h + u)}} - \frac{\bar{c}}{\sqrt{(h + \bar{u})}}\right\};
\]
which are in fact the equations which express that $Q$ and $\bar{Q}$ are corresponding points.

\medskip

It is proper to remark that supposing, as we are at liberty to do, that $P$, $\bar{P}$ are situate in corresponding octants of the two ellipsoids, then if the curve of contact of the circumscribed cone having $P$ for its vertex divide the surface of the interior ellipsoid into two parts $\bar{M}$, $\bar{N}$, of which the former lies contiguous to $\bar{P}$: also if the curve of intersection of the tangent plane at $\bar{P}$ divide the surface of the exterior ellipsoid into two parts $M$, $N$, of which $M$ lies contiguous to the point $P$; then the different points of $M$, $\bar{M}$ correspond to each other, as do also the different points of $N$, $\bar{N}$.

\medskip

We may now pass to the integral calculus problem. The Attraction parallel to the axis of $x$ is
\[
A = \int \frac{x\, dx\, dy\, dz}{(x^{2} + y^{2} + z^{2})^{\frac{3}{2}}}\,,
\]
the limits of the integration being given $(\bar{a}^{2})$
\[
\frac{x^{2}}{f + \bar{u}} + \frac{y^{2}}{g + \bar{u}} + \frac{z^{2}}{h + \bar{u}} = 1\,;
\]
or putting
\[
x = r X,\qquad y = r Y,\qquad z = r Z,
\]
where $X$, $Y$, $Z$ have the same signification as before, we have
\[
dx\, dy\, dz = r^{2}\, dr\, dS,
\]
and then
\[
A = \int X\, dr\, dS = \int \rho X\, dS,
\]
where $\rho$ has the same signification as before: it will be convenient to leave the formula in this form, rather than to take at once the difference of the two values of $\rho$, but of course the integration is as in the ordinary methods to be performed so as to extend to the whole volume of the ellipsoid. The expression $dS$ denotes the differen-

\medskip
\noindent\textit{[Book p.~516]}\hfill

\smallskip

\noindent tial of a spherical surface radius unity, and if $\theta$, $\phi$ are the parameters by which the position of $\rho$ is determined, we have
\[
dS = \left| \begin{array}{ccc} X, & Y, & Z \\ \dfrac{dX}{d\theta}, & \dfrac{dY}{d\theta}, & \dfrac{dZ}{d\theta} \\ \dfrac{dX}{d\phi}, & \dfrac{dY}{d\phi}, & \dfrac{dZ}{d\phi} \end{array} \right| d\theta\, d\phi.
\]
In the present case
\[
dS = dS_{1} = \left| \begin{array}{ccc} X_{1}, & Y_{1}, & Z_{1} \\ \dfrac{dX_{1}}{dY_{1}}, & \dfrac{dY_{1}}{dY_{1}}, & \dfrac{dZ_{1}}{dY_{1}} \\ \dfrac{dX_{1}}{dZ_{1}}, & \dfrac{dY_{1}}{dZ_{1}}, & \dfrac{dZ_{1}}{dZ_{1}} \end{array} \right| dY_{1}\, dZ_{1}.
\]
or from the values of $X_{1}$, $Y_{1}$, $Z_{1}$ in terms of $\bar{X}_{1}$, $\bar{Y}_{1}$, $\bar{Z}_{1}$ (observing that $X_{1}$ must be replaced by its value $\sqrt{(1 - Y_{1}^{2} - Z_{1}^{2})}$) we deduce
\[
dS = \sqrt{\!\left[\frac{(\bar{u} - q)(\bar{u} - r)}{(u - q)(u - r)}\right]}\,\frac{1}{X_{1}}\, d\bar{Y}_{1}\, d\bar{Z}_{1}.
\]
But $d\bar{S} = d\bar{S}_{1} = \dfrac{1}{\bar{X}_{1}}\, d\bar{Y}_{1}\, d\bar{Z}_{1}$, whence
\[
dS = \sqrt{\!\left[\frac{(\bar{u} - q)(\bar{u} - r)}{(u - q)(u - r)}\right]}\,\frac{\bar{X}_{1}\, d\bar{S}}{X_{1}}\,,
\]
which shows that the corresponding elements of the spheres whose centres are $P$, $\bar{P}$, projected upon the tangent planes at $P$ and $\bar{P}$ respectively, are in a constant ratio. It may be noticed also that if $\mu$, $\bar{\mu}$ are the masses of the ellipsoids, the ratio in question
\[
= \sqrt{\!\left[\frac{(\bar{u} - q)(\bar{u} - r)}{(u - q)(u - r)}\right]} = \frac{\bar{\mu} \bar{a}_{1}}{\mu \bar{a}_{1}}\,.
\]
We have thus
\[
A = \sqrt{\!\left[\frac{(\bar{u} - q)(\bar{u} - r)}{(u - q)(u - r)}\right]} \int \frac{\bar{\rho}\bar{X}_{1} X\, d\bar{S}}{X_{1}}\,,
\]
that is
\[
A = \sqrt{\!\left[\frac{(\bar{u} - q)(\bar{u} - r)}{(u - q)(u - r)}\right]} \int \frac{\bar{\rho}\,(\alpha X_{1} + \alpha' Y_{1} X_{1} + \alpha'' Z_{1})\, X_{1}\, d\bar{S}}{X_{1}}\,.
\]

The value which it will be convenient to use for $\rho$ is that derived from the equation
\[
\bar{\rho}^{2}\!\left(\frac{X^{2}}{f + \bar{u}} + \frac{Y^{2}}{g + \bar{u}} + \frac{Z^{2}}{h + \bar{u}}\right) + 2 \bar{\rho}\!\left(\frac{\bar{X}\bar{a}}{f + \bar{u}} + \frac{\bar{Y}\bar{b}}{g + \bar{u}} + \frac{\bar{Z}\bar{c}}{h + \bar{u}}\right) + \frac{\bar{u} - u}{\bar{a}_{1}^{2}} = 0.
\]

\medskip
\noindent\textit{[Book p.~517]}\hfill

\smallskip

\noindent with only the transformation of expressing the radical in terms of $X_{1}$, viz.\
\[
\bar{\rho} = \frac{-\dfrac{\bar{X}\bar{a}}{f + \bar{u}} - \dfrac{\bar{Y}\bar{b}}{g + \bar{u}} - \dfrac{\bar{Z}\bar{c}}{h + \bar{u}} + \dfrac{1}{\bar{a}_{1}}\, X_{1}}{\dfrac{X^{2}}{f + \bar{u}} + \dfrac{Y^{2}}{g + \bar{u}} + \dfrac{Z^{2}}{h + \bar{u}}}\,;
\]
substituting these values and observing that $Y_{1}$ and $Z_{1}$ are rational functions of $X$, $Y$, and $Z$, but that $X_{1}$ is a radical, and that in order to extend the integration to the whole ellipsoid, the values corresponding to the opposite signs of $X_{1}$ will require to be added, the quantity to be integrated (omitting for the moment the exterior constant factor) is
\[
\dfrac{-\!\left(\dfrac{\bar{X}\bar{a}}{f + \bar{u}} + \dfrac{\bar{Y}\bar{b}}{g + \bar{u}} + \dfrac{\bar{Z}\bar{c}}{h + \bar{u}}\right) \dfrac{1}{\bar{a}_{1}}\,\{\alpha' Y_{1} + \alpha'' Z_{1}\}\, X_{1}\, d\bar{S}}{\dfrac{X^{2}}{f + \bar{u}} + \dfrac{Y^{2}}{g + \bar{u}} + \dfrac{Z^{2}}{h + \bar{u}}}\,,
\]
the integration to be extended over the spherical area $\bar{S}$. Consider the quantity within $\{\,\}$, this is
\[ \resizebox{\textwidth}{!}{$\displaystyle
\alpha\!\left(\dfrac{\bar{X}\bar{a}}{f + \bar{u}} + \dfrac{\bar{Y}\bar{b}}{g + \bar{u}} + \dfrac{\bar{Z}\bar{c}}{h + \bar{u}}\right) + \dfrac{\alpha'}{\bar{a}_{1}}\sqrt{\!\left(\dfrac{\bar{u} - q}{u - q}\right)}\!\left(\bar{\alpha}'\bar{X} + \bar{\beta}'\bar{Y} + \bar{\gamma}'\bar{Z}\right) + \dfrac{\alpha''}{\bar{a}_{1}}\sqrt{\!\left(\dfrac{\bar{u} - r}{u - r}\right)}\!\left(\bar{\alpha}''\bar{X} + \bar{\beta}''\bar{Y} + \bar{\gamma}''\bar{Z}\right).
$} \]

The coefficients of $\bar{Y}$ and $\bar{Z}$ vanish, in fact that of $\bar{Y}$ is
\[ \resizebox{\textwidth}{!}{$\displaystyle
\frac{a_{1} a}{f + u}\,\frac{\bar{b}}{g + \bar{u}} + \frac{b_{1} a}{a_{1}(f + q)}\sqrt{\!\left(\frac{\bar{u} - q}{u - q}\right)}\,\frac{\bar{b}_{1} \bar{b}}{g + q} + \frac{c_{1} a}{a_{1}(f + r)}\sqrt{\!\left(\frac{\bar{u} - r}{u - r}\right)}\,\frac{\bar{c}_{1} \bar{b}}{g + r}
$} \]
\[ \resizebox{\textwidth}{!}{$\displaystyle
= \frac{a \bar{b}}{a_{1}}\!\left\{\frac{a_{1}^{2}}{(f + u)(g + \bar{u})} + \frac{b_{1} \bar{b}_{1}}{(f + q)(g + q)}\sqrt{\!\left(\frac{\bar{u} - q}{u - q}\right)} + \frac{c_{1} \bar{c}_{1}}{(f + r)(g + r)}\sqrt{\!\left(\frac{\bar{u} - r}{u - r}\right)}\right\}
$} \]
\[
= \frac{a \bar{b}}{a_{1}}\!\left\{\frac{a_{1}^{2}}{(f + u)(g + u)} + \frac{b_{1}^{2}}{(f + q)(g + q)} + \frac{c_{1}^{2}}{(f + r)(g + r)}\right\} = 0,
\]
and similarly for the coefficient of $\bar{Z}$.

The coefficient of $\bar{X}$ is in like manner shown to be
\[
\frac{a\bar{a}}{a_{1}}\!\left\{\frac{a_{1}^{2}}{(f + u)^{2}} + \frac{b_{1}^{2}}{(f + q)^{2}} + \frac{c_{1}^{2}}{(f + r)^{2}}\right\} = \frac{a\bar{a}}{a_{1}}\cdot\frac{(f - g)(f - h)}{(f + u)(f + q)(f + r)} = \frac{\bar{a}}{a^{2} a_{1}} = \frac{\bar{a}}{a a_{1}'}\,,
\]
or the quantity in question is simply
\[
\frac{\bar{a}}{a a_{1}}\,\bar{X}.
\]

\medskip
\noindent\textit{[Book p.~518]}\hfill

\smallskip

\noindent Multiplying this by $\bar{X}_{1} = \bar{\alpha}\bar{X} + \bar{\alpha}'\bar{Y} + \bar{\alpha}''\bar{Z}$, the terms containing $\bar{X}\bar{Y}$, $\bar{X}\bar{Z}$ vanish after the integration, and we need only consider the term $\bar{\alpha}\bar{X}^{2}$, or what is the same $\dfrac{\bar{a}^{2} \bar{a}_{1}}{a a_{1}(f + \bar{u})}\,\bar{X}^{2}$. Hence
\[
A = \sqrt{\!\left[\frac{(\bar{u} - q)(\bar{u} - r)}{(u - q)(u - r)}\right]}\,\frac{\bar{a}^{2} \bar{a}_{1}}{a a_{1}(f + \bar{u})}\!\int \frac{\bar{X}^{2}\, d\bar{S}}{\dfrac{\bar{X}^{2}}{f + \bar{u}} + \dfrac{\bar{Y}^{2}}{g + \bar{u}} + \dfrac{\bar{Z}^{2}}{h + \bar{u}}}\,.
\]

The value of the corresponding function $\bar{A}$ (that is, the attraction of the exterior ellipsoid upon $\bar{P}$) is
\[
\bar{A} = \frac{\bar{a}}{f + \bar{u}}\!\int \frac{\bar{X}^{2}\, d\bar{S}}{\dfrac{\bar{X}^{2}}{f + \bar{u}} + \dfrac{\bar{Y}^{2}}{g + \bar{u}} + \dfrac{\bar{Z}^{2}}{h + \bar{u}}}\,,
\]
the limits being the same, whence
\[
\frac{A}{\bar{A}} = \sqrt{\!\left[\frac{(\bar{u} - q)(\bar{u} - r)}{(u - q)(u - r)}\right]}\,\frac{f + \bar{u}}{f + \bar{u}}\cdot\frac{\bar{a} a_{1}}{a \bar{a}_{1}} = \frac{\sqrt{(u + g)}\,\sqrt{(u + h)}}{\sqrt{(\bar{u} + g)}\,\sqrt{(\bar{u} + h)}}\,,
\]
or we have
\[
A = \frac{\sqrt{(u + g)}\,\sqrt{(u + h)}}{\sqrt{(\bar{u} + g)}\,\sqrt{(\bar{u} + h)}}\,\bar{A},\qquad B = \frac{\sqrt{(u + h)}\,\sqrt{(u + f)}}{\sqrt{(\bar{u} + h)}\,\sqrt{(\bar{u} + f)}}\,\bar{B},\qquad C = \frac{\sqrt{(u + f)}\,\sqrt{(u + g)}}{\sqrt{(\bar{u} + f)}\,\sqrt{(\bar{u} + g)}}\,\bar{C},
\]
formulae which constitute in fact Ivory's theorem.

Let $K$, $\bar{K}$ denote the attractions in the directions of the normals at $P$, $\bar{P}$, we have
\[
K = \int X\, d\bar{S},\qquad \bar{K} = \int \bar{X}\, d\bar{S},
\]
or
\[
K = \frac{\bar{\mu} a_{1}}{\mu \bar{a}_{1}}\,\bar{K}\,;
\]
and it is important to remark that this is true not only for the entire ellipsoids; but if $\mathfrak{M}$, $\mathfrak{N}$ denote the attractions of the cones standing on the portions $\bar{M}$, $\bar{N}$ of the surface of the interior ellipsoid, and $\bar{\mathfrak{M}}$, $\bar{\mathfrak{N}}$ the attractions of the portions of the exterior ellipsoid bounded by the tangent plane at $\bar{P}$, and the portions $M$, $N$ of the surface of the exterior ellipsoid, then
\[
\mathfrak{M} = \frac{\bar{\mu} a_{1}}{\mu \bar{a}_{1}}\,\bar{\mathfrak{M}},\qquad \mathfrak{N} = \frac{\bar{\mu} a_{1}}{\mu \bar{a}_{1}}\,\bar{\mathfrak{N}},
\]
where obviously
\[
K = \mathfrak{M} - \mathfrak{N},\qquad \bar{K} = \bar{\mathfrak{M}} + \bar{\mathfrak{N}}\,;
\]
this theorem is so far as I am aware new.

\bigskip
\hrule
\bigskip

% ===================================================================
% Paper [90] -- NOTE SUR QUELQUES FORMULES RELATIVES AUX CONIQUES.
% Book pp. 519--521.  (French, verbatim.)
% ===================================================================

\section*{[90] Note sur quelques formules relatives aux coniques}

\noindent\textit{[Book p.~519]}\hfill

\smallskip

\noindent\emph{[From the Journal f\"ur die reine und angewandte Mathematik (Crelle), tom.~\textsc{xxxix}.\ (1850), pp.~1--3.]}

\medskip

\noindent Soit, comme \`a l'ordinaire:
\[
\begin{aligned}
\mathfrak{A} &= B C - F^{2},\\
\mathfrak{B} &= C A - G^{2},\\
\mathfrak{C} &= A B - H^{2},\\
\mathfrak{F} &= G H - A F,\\
\mathfrak{G} &= H F - B G,\\
\mathfrak{H} &= F G - C H,\\
K &= A B C - A F^{2} - B G^{2} - C H^{2} + 2 F G H,
\end{aligned}\qquad\quad\cdots\cdots (1)
\]
et d\'esignons, pour abr\'eger, les fonctions
\[
\begin{aligned}
& A x^{2} + B y^{2} + C z^{2} + 2 F y z + 2 G x z + 2 H x y,\\
& A \alpha x + B \beta y + C \gamma z + F(\beta z + \gamma y) + G(\gamma x + \alpha z) + H(\alpha y + \beta x),\\
& A(\beta z - \gamma y)^{2} + B(\gamma x - \alpha z)^{2} + C(\alpha y - \beta x)^{2} + 2 F(\gamma x - \alpha z)(\alpha y - \beta x)\\
& \qquad\qquad + 2 G(\alpha y - \beta x)(\beta z - \gamma y) + 2 H(\beta z - \gamma y)(\gamma x - \alpha z);\\
& \&c.
\end{aligned}
\]
par
\[
A x^{2} + \dots\,;\quad A \alpha x + \dots\,;\quad A(\beta z - \gamma y)^{2} + \dots\,;\ \&c.
\]

Cela pos\'e, soient $\mathfrak{A} + \mathfrak{a}$, $\mathfrak{B} + \mathfrak{b}$, $\mathfrak{C} + \mathfrak{c}$, $\mathfrak{F} + \mathfrak{f}$, $\mathfrak{G} + \mathfrak{g}$, $\mathfrak{H} + \mathfrak{h}$, $K + k$ ce que deviennent $\mathfrak{A}$, $\mathfrak{B}$, $\mathfrak{C}$, $\mathfrak{F}$, $\mathfrak{G}$, $\mathfrak{H}$, $K$, en \'ecrivant $A + \alpha^{2}$, $B + \beta^{2}$, $C + \gamma^{2}$, $F + \beta\gamma$, $G + \gamma\alpha$, $H + \alpha\beta$ au lieu de $A$, $B$, $C$, $F$, $G$, $H$. Alors on aura d'abord
\[
k = \mathfrak{A}\alpha^{2} + \dots \qquad\qquad\cdots\cdots\cdots\cdots\cdots\cdots\cdots\cdots\cdots\cdots (2)
\]

\medskip
\noindent\textit{[Book p.~520]}\hfill

\smallskip

\noindent et les quantit\'es $\mathfrak{a}$, $\mathfrak{b}$, $\mathfrak{c}$, $\mathfrak{f}$, $\mathfrak{g}$, $\mathfrak{h}$ seront donn\'ees par l'\'equation
\[
\mathfrak{a} \alpha^{2} + \dots = A(\beta z - \gamma y)^{2} + \dots \qquad\qquad\cdots\cdots\cdots\cdots (3),
\]
ou, si l'on veut, par celle-ci:
\[
\mathfrak{a} \alpha x_{1} + \dots = A(\beta z - \gamma y)(\beta z_{1} - \gamma y_{1}) + \dots \qquad\cdots\cdots\cdots (4),
\]
(savoir, en consid\'erant ces \'equations comme identiques par rapport \`a $x$, $y$, $z$ et $x_{1}$, $y_{1}$, $z_{1}$, respectivement).

On obtient sans difficult\'e les \'equations identiques:
\[
(A \alpha x_{1} + \dots)(A x x_{1} + \dots) - (A \alpha x + \dots)(A \alpha_{1} x_{1} + \dots) = \mathfrak{A}(\beta z_{1} - \gamma y_{1})(\beta z - \gamma y) + \dots \quad (5),
\]
\[
(\mathfrak{a} \alpha \alpha_{1} + \dots)(\mathfrak{a} x x_{1} + \dots) - (\mathfrak{a} \alpha x + \dots)(\mathfrak{a} \alpha_{1} x_{1} + \dots) = -K\,[\,A(\beta z_{1} - \gamma y_{1})(\beta z - \gamma y) + \dots\,] \quad (6).
\]

Comme on sait, la condition sous laquelle la droite $l x + m y + n z = 0$ touche la conique $U = A x^{2} + \dots = 0$ peut \^etre pr\'esent\'ee sous la forme $\mathfrak{A} l^{2} + \dots = 0$. Donc la condition pour que cette droite touche la conique $U + (\alpha x + \beta y + \gamma z)^{2} = 0$, est
\[
\mathfrak{A} l^{2} + \dots + A(\gamma m - \beta n)^{2} + \dots = 0 \qquad\cdots\cdots\cdots\cdots (7).
\]
En r\'eduisant au moyen de l'\'equation $(\mathfrak{A}\alpha^{2} + \dots)(\mathfrak{A} l^{2} + \dots) = K\,[\,A(\gamma m - \beta n)^{2} + \dots\,]$ (laquelle n'est qu'un cas particulier de l'\'equation (6)), cette condition devient:
\[
(K + \mathfrak{A}\alpha^{2} + \dots)(\mathfrak{A} l^{2} + \dots) - (\mathfrak{A}\alpha l + \dots)^{2} = 0 \qquad\cdots\cdots\cdots (8).
\]

Pour trouver la condition sous laquelle les coniques
\[
U + (\alpha x + \beta y + \gamma z)^{2} = 0,\qquad U + (\alpha_{1} x + \beta_{1} y + \gamma_{1} z)^{2} = 0
\]
se touchent, on n'a qu'\`a remarquer que l'\'equation de la tangente commune est, ou
\[
(\alpha - \alpha_{1}) x + (\beta - \beta_{1}) y + (\gamma - \gamma_{1}) z = 0,\quad \text{ou}\quad (\alpha + \alpha_{1}) x + (\beta + \beta_{1}) y + (\gamma + \gamma_{1}) z = 0.
\]
En ne consid\'erant que la premi\`ere de ces droites, on a pour la condition cherch\'ee:
\[
\mathfrak{A}(\alpha - \alpha_{1})^{2} + \dots + A(\beta\gamma_{1} - \beta_{1}\gamma)^{2} + \dots = 0 \qquad\cdots\cdots\cdots\cdots (9);
\]
ou, en r\'eduisant au moyen du m\^eme cas particulier de l'\'equation (6), on obtient cette condition sous la forme
\[
(K + \mathfrak{A}\alpha^{2} + \dots)(K + \mathfrak{A}\alpha_{1}^{2} + \dots) - (K + \mathfrak{A}\alpha\alpha_{1} + \dots)^{2} = 0 \qquad\cdots\cdots (10).
\]

On sait que les coordonn\'ees $X$, $Y$, $Z$ du p\^ole de la droite $l x + m y + n z = 0$, par rapport \`a la conique $U = 0$, peuvent \^etre trouv\'ees par l'\'equation
\[
\lambda X + \mu Y + \nu Z = \mathfrak{A} l \lambda + \dots,
\]

\medskip
\noindent\textit{[Book p.~521]}\hfill

\smallskip

\noindent (consid\'er\'ee comme identique par rapport \`a $\lambda$, $\mu$, $\nu$). De l\`a les coordonn\'ees $X$, $Y$, $Z$ du p\^ole de cette m\^eme droite par rapport \`a la conique
\[
U + (\alpha x + \beta y + \gamma z)^{2} = 0
\]
se trouveront par l'\'equation
\[
\lambda X + \mu Y + \nu Z = \mathfrak{A} l \lambda + \dots + A(\gamma m - \beta n)(\gamma\mu - \beta\nu). \qquad (11)
\]
(consid\'er\'ee comme identique par rapport \`a $\lambda$, $\mu$, $\nu$). Cette \'equation peut aussi \^etre pr\'esent\'ee sous la forme
\[
K(\lambda X + \mu Y + \nu Z) = (K + \mathfrak{A}\alpha^{2} + \dots)(\mathfrak{A} l \lambda + \dots) - (\mathfrak{A}\alpha\lambda + \dots)(\mathfrak{A}\alpha l + \dots) \quad \cdots (12),
\]
ce qui peut \^etre d\'emontr\'e facilement au moyen d'un cas particulier de l'\'equation (6).

Si les deux droites $l x + m y + n z = 0$, $l' x + m' y + n' z = 0$, touchent la conique $U = 0$, l'\'equation de la droite qui passe par les points de contact sera $A(m n' - m' n) x + \dots = 0$. Donc: si ces deux droites touchent la conique $U + (\alpha x + \beta y + \gamma z)^{2} = 0$, l'\'equation de la droite qui passe par les deux points de contact sera
\[ \resizebox{\textwidth}{!}{$\displaystyle
A(m n' - m' n)\,x + \dots + (\alpha x + \beta y + \gamma z)\,[\,\alpha(m n' - m' n) + \beta(n l' - n' l) + \gamma(l m' - l' m)\,] = 0 \quad \cdots (13).
$} \]

Les formules obtenues seront utiles pour la solution du probl\`eme du m\'emoire suivant, [91]. Je les ai rapproch\'ees ici pour ne pas interrompre cette solution.

\bigskip
\hrule
\bigskip

% ===================================================================
% Paper [91] -- SUR LE PROBLEME DES CONTACTS.
% Book pp. 522--523.  (French, verbatim, only opening pages in this range.)
% ===================================================================

\section*{[91] Sur le probl\`eme des contacts}

\noindent\textit{[Book p.~522]}\hfill

\smallskip

\noindent\emph{[From the Journal f\"ur die reine und angewandte Mathematik (Crelle), tom.~\textsc{xxxix}.\ (1850), pp.~4--13.]}

\medskip

\noindent Je me propose ici la solution analytique du probl\`eme suivant:

``\'Etant donn\'ees trois coniques inscrites \`a une m\^eme conique: trouver une autre conique, aussi inscrite \`a cette conique, qui touche les trois coniques inscrites; et tirer de l\`a les constructions g\'eom\'etriques ordinaires.''

\medskip

Je commence par r\'ecapituler quelques-unes des propri\'et\'es d'un syst\`eme de trois coniques inscrites \`a la m\^eme conique.

\medskip

Un syst\`eme de six droites qui passent trois \`a trois par quatre points, s'appelle \emph{quadrangle}. Les points de rencontre des c\^ot\'es oppos\'es sont les \emph{centres} du quadrangle; les c\^ot\'es du triangle form\'e par ces trois centres sont les \emph{axes} du quadrangle. De m\^eme, un syst\`eme de six points situ\'es trois \`a trois sur quatre droites, s'appelle \emph{quadrilat\`ere}. Les droites qui passent par les angles oppos\'es sont les \emph{axes}; et les angles du triangle form\'e par les trois axes sont les \emph{centres} du quadrilat\`ere.

Deux coniques quelconques se coupent en quatre points qui forment un \emph{quadrangle} inscrit aux deux coniques. Elles ont quatre tangentes communes qui forment un \emph{quadrilat\`ere} circonscrit aux deux coniques. Le quadrangle inscrit et le quadrilat\`ere circonscrit ont les m\^emes centres et les m\^emes axes.

Si deux coniques sont circonscrites ou inscrites l'une \`a l'autre, la droite qui passe par les deux points de contact s'appelle \emph{chorde de contact}, et le point de rencontre des deux tangentes communes s'appelle \emph{centre de contact}.

Cela pos\'e: les coniques circonscrites \`a deux coniques donn\'ees, peuvent \^etre divis\'ees en trois classes: une conique circonscrite appartient \`a une quelconque de ces trois classes, selon que les points de rencontre des chordes de contact de la conique

\medskip
\noindent\textit{[Book p.~523]}\hfill

\smallskip

\noindent circonscrite et de chacune des deux coniques donn\'ees coincide avec un quelconque des trois centres du quadrangle inscrit, ou du quadrilat\`ere circonscrit; ou, si l'on veut, selon que la droite qui passe par les centres de contact de la conique circonscrite et des deux coniques donn\'ees, coincide avec un quelconque des trois axes du quadrangle inscrit, ou du quadrilat\`ere circonscrit.

En consid\'erant les deux coniques donn\'ees, et une conique circonscrite, nous dirons que les deux c\^ot\'es du quadrangle inscrit qui se coupent dans le point d'intersection des deux chordes de contact, sont les \emph{axes de symptose} des deux coniques donn\'ees, et que les deux angles du quadrilat\`ere circonscrit, situ\'es sur la droite qui passe par les deux centres de contact, sont les \emph{centres d'homologie} des deux coniques donn\'ees.

Soient maintenant inscrites trois coniques \`a la m\^eme conique. En combinant deux \`a deux ces trois coniques, les six axes de symptose se couperont trois \`a trois en quatre points que nous nommerons \emph{centres de symptose}, et les six centres d'homologie seront situ\'es trois \`a trois sur quatre droites que nous appelerons \emph{axes d'homologie}.

Soient
\[
U + V_{1}^{2} = 0,\qquad U + V_{2}^{2} = 0,\qquad U + V_{3}^{2} = 0
\]
les \'equations des trois coniques inscrites, ou
\[
\begin{aligned}
U &= A x^{2} + B y^{2} + C z^{2} + 2 F y z + 2 G x z + 2 H x y,\\
V_{1} &= \alpha_{1} x + \beta_{1} y + \gamma_{1} z,\\
V_{2} &= \alpha_{2} x + \beta_{2} y + \gamma_{2} z,\\
V_{3} &= \alpha_{3} x + \beta_{3} y + \gamma_{3} z.
\end{aligned}
\]

Si
\[
l x + m y + n z = 0
\]
est l'\'equation d'une tangente commune aux coniques $U + V_{1}^{2} = 0$, $U + V_{2}^{2} = 0$, les formules de la ``Note sur quelques formules \&c.'' [90], en adoptant la notation de cette note, donneront les \'equations
\[
\begin{aligned}
(K + \mathfrak{A}\alpha_{1}^{2} + \dots)(\mathfrak{A} l^{2} + \dots) - (\mathfrak{A}\alpha_{1} l + \dots)^{2} &= 0,\\
(K + \mathfrak{A}\alpha_{2}^{2} + \dots)(\mathfrak{A} l^{2} + \dots) - (\mathfrak{A}\alpha_{2} l + \dots)^{2} &= 0,
\end{aligned}
\]
qui serviront \`a d\'eterminer les valeurs de $l$, $m$, $n$; on obtient par l\`a l'expression
\[
\sqrt{(K + \mathfrak{A}\alpha_{1}^{2} + \dots)}\,(\mathfrak{A}\alpha_{2} l + \dots) - \sqrt{(K + \mathfrak{A}\alpha_{2}^{2} + \dots)}\,(\mathfrak{A}\alpha_{1} l + \dots),
\]
qui fait voir que l'\'equation $l x + m y + n z = 0$ est satisfaite en \'ecrivant
\[ \resizebox{\textwidth}{!}{$\displaystyle
\sqrt{(K + \mathfrak{A}\alpha_{1}^{2} + \dots)}\,(\mathfrak{G}\alpha_{2} + \mathfrak{F}\beta_{2} + \mathfrak{C}\gamma_{2}) - \sqrt{(K + \mathfrak{A}\alpha_{2}^{2} + \dots)}\,(\mathfrak{G}\alpha_{1} + \mathfrak{F}\beta_{1} + \mathfrak{C}\gamma_{1}).
$} \]

\end{document}
